% ==========================================================================
%  Figure 1: DeepLoop overview.
%  Self-contained: include with % ==========================================================================
%  Figure 1: DeepLoop overview.
%  Self-contained: include with % ==========================================================================
%  Figure 1: DeepLoop overview.
%  Self-contained: include with % ==========================================================================
%  Figure 1: DeepLoop overview.
%  Self-contained: include with \input{figures/fig1_overview.tex}.
%  Requires: tikz (+ arrows.meta, calc libraries), amsmath.
%  Helvetica (phv, scaled 0.92) is scoped to this group only.
% ==========================================================================
\begingroup
\makeatletter
\def\Hv@scale{0.92}%
\makeatother
\renewcommand{\sfdefault}{phv}%
% ---- palette: two colorblind-safe accents (block identity) + neutrals ----
\definecolor{dlBlue}  {HTML}{4C72B0} % Block 1 / phi_1
\definecolor{dlOrange}{HTML}{DD8452} % Block 2 / phi_2
\definecolor{dlInk}   {HTML}{2B2F36} % main text
\definecolor{dlMut}   {HTML}{6A7280} % muted annotations
\definecolor{dlLine}  {HTML}{9AA0A8} % neutral borders
\definecolor{dlArrow} {HTML}{565D66} % arrows
\begin{tikzpicture}[
  font=\sffamily,
  >={Stealth[length=1.7mm,width=1.4mm]},
  line cap=round, line join=round,
  every node/.style={text=dlInk, outer sep=0pt},
  % small sublayer box (unrolled chain)
  subS/.style 2 args={rectangle, rounded corners=2pt, draw=#1,
      line width=0.6pt, fill=#2, minimum width=7.6mm, minimum height=6.2mm,
      inner sep=1pt, font=\sffamily\scriptsize},
  % large sublayer box (physical panel)
  subL/.style 2 args={rectangle, rounded corners=2pt, draw=#1,
      line width=0.6pt, fill=#2, minimum width=11.5mm, minimum height=6.8mm,
      inner sep=1pt, font=\sffamily\footnotesize},
  % state chips x_0, x_N, x_i, ...
  chip/.style={rectangle, rounded corners=2pt, draw=dlLine, fill=black!5,
      line width=0.6pt, minimum height=6.2mm, minimum width=8.5mm,
      inner sep=2.5pt, font=\footnotesize},
  % neutral gray box (Norm, f_j)
  gray box/.style={rectangle, rounded corners=2pt, draw=dlLine, fill=black!4,
      line width=0.6pt, font=\sffamily\footnotesize},
  flow/.style={->, draw=dlArrow, line width=0.6pt,
      shorten >=1.2pt, shorten <=1.2pt},
  head/.style={font=\sffamily\small, anchor=base west, inner sep=0pt},
  mut/.style={font=\sffamily\scriptsize, text=dlMut},
]

% =========================================================================
%  (a) PHYSICAL BLOCKS
% =========================================================================
\node[head] at (-1.53, 10.15)
  {{\bfseries (a)\; Physical blocks} \textcolor{dlMut}{---\;
   $K{=}2$ shared blocks, stored once}};

% Block 1
\node[rectangle, rounded corners=2pt, draw=dlBlue, fill=dlBlue!6,
      line width=0.8pt, minimum width=38mm, minimum height=16mm]
  (B1) at (0.77, 8.90) {};
\node[font=\sffamily\footnotesize\bfseries, text=dlBlue]
  at (0.77, 9.37) {Block 1\; ($\phi_1$)};
\node[subL={dlBlue}{dlBlue!13}, minimum width=13mm]
  (b1a) at (-0.08, 8.60) {attn};
\node[subL={dlBlue}{dlBlue!30}, minimum width=13mm]
  (b1f) at ( 1.62, 8.60) {ffn};
\draw[flow] (b1a) -- (b1f);

% Block 2
\node[rectangle, rounded corners=2pt, draw=dlOrange, fill=dlOrange!6,
      line width=0.8pt, minimum width=38mm, minimum height=16mm]
  (B2) at (5.37, 8.90) {};
\node[font=\sffamily\footnotesize\bfseries, text=dlOrange]
  at (5.37, 9.37) {Block 2\; ($\phi_2$)};
\node[subL={dlOrange}{dlOrange!13}, minimum width=13mm]
  (b2a) at (4.52, 8.60) {attn};
\node[subL={dlOrange}{dlOrange!30}, minimum width=13mm]
  (b2f) at (6.22, 8.60) {ffn};
\draw[flow] (b2a) -- (b2f);

\draw[flow] (B1) -- (B2);

% loop-back arc under the two blocks, label set on the wire
\draw[flow, rounded corners=3pt]
  (B2.east) -- ++(0.40,0) |- (3.07, 7.55) -| ($(B1.west)+(-0.40,0)$)
  -- (B1.west);
\node[font=\sffamily\footnotesize, fill=white, inner sep=2.5pt]
  at (3.07, 7.55) {loop $\times\,R$ rounds};

% =========================================================================
%  (b) UNROLLED EXECUTION
% =========================================================================
\node[head] at (-1.53, 6.45)
  {{\bfseries (b)\; Unrolled execution} \textcolor{dlMut}{---\;
   the same two blocks are revisited in every round}};

\def\yc{5.45}
% input / output chips
\node[chip] (xin)  at (-1.105, \yc) {$\mathbf{x}_0$};
\node[chip] (xout) at (14.085, \yc) {$\mathbf{x}_N$};

% six unrolled 2-sublayer groups
\foreach \gx/\hue/\pj [count=\i] in
  {0/dlBlue/1, 2.30/dlOrange/2, 4.82/dlBlue/1,
   7.12/dlOrange/2, 9.64/dlBlue/1, 11.94/dlOrange/2}{
  \node[subS={\hue}{\hue!13}] (a\i) at (\gx, \yc) {attn};
  \node[subS={\hue}{\hue!30}] (f\i) at (\gx+1.04, \yc) {ffn};
  \draw[flow] (a\i) -- (f\i);
  \draw[draw=\hue!45, line width=0.5pt, rounded corners=2pt]
    (\gx-0.51, \yc-0.45) rectangle (\gx+1.55, \yc+0.45);
  \node[font=\sffamily\scriptsize, text=\hue] at (\gx+0.52, \yc-0.72)
    {$\phi_{\pj}$};
}

% chain arrows
\draw[flow] (xin) -- (a1);
\foreach \i [evaluate=\i as \j using int(\i+1)] in {1,...,5}{
  \draw[flow] (f\i) -- (a\j);}
\draw[flow] (f6) -- (xout);

% round dividers + labels
\draw[dlLine!70, line width=0.45pt, dash pattern=on 2.6pt off 2pt]
  (4.08, \yc+0.5) -- (4.08, \yc-1.25);
\draw[dlLine!70, line width=0.45pt, dash pattern=on 2.6pt off 2pt]
  (8.90, \yc+0.5) -- (8.90, \yc-1.25);
\node[mut] at ( 1.67, \yc-1.05) {round $r{=}1$};
\node[mut] at ( 6.49, \yc-1.05) {round $r{=}2$};
\node[mut] at (11.31, \yc-1.05) {round $r{=}3$};

% summary line
\node[font=\sffamily\footnotesize, text=dlInk] at (6.50, 3.90)
  {$R{=}3$ rounds $\times$ $K{=}2$ blocks $\;\Rightarrow\;$
   unrolled depth $N{=}KR{=}6$,\quad $M{=}2N{=}12$ sublayer visits};

% =========================================================================
%  (c) RESIDUAL SUBLAYER
% =========================================================================
\node[head] at (-1.53, 3.10) {{\bfseries (c)\; Residual sublayer:}\;\;
  $\mathbf{x}_{i+1} = \mathrm{Norm}\bigl(\alpha\,\mathbf{x}_i
    + f_j(\mathrm{Norm}(\mathbf{x}_i);\phi_j)\bigr)$};

% dataflow: skip path is the straight wire, branch dips through f_j
\def\yd{1.95}
\node[chip] (ri) at (-1.105, \yd) {$\mathbf{x}_i$};
\node[circle, draw=dlArrow, fill=white, line width=0.6pt,
      minimum size=4.2mm, inner sep=0pt, font=\footnotesize]
  (add) at (5.05, \yd) {$+$};
\node[gray box, minimum width=13.5mm, minimum height=6.5mm]
  (nrm) at (6.90, \yd) {Norm};
\node[chip] (ro) at (8.80, \yd) {$\mathbf{x}_{i+1}$};
\node[gray box, minimum width=10mm, minimum height=6.5mm]
  (inrm) at (1.15, 0.85) {Norm};
\node[gray box, minimum width=21mm, minimum height=8mm]
  (fj) at (3.10, 0.85) {$f_j(\,\cdot\,;\phi_j)$};

% split point
\coordinate (sp) at (0.25, \yd);
\fill[dlArrow] (sp) circle (1.1pt);
\draw[draw=dlArrow, line width=0.6pt] (ri) -- (sp);
% skip path (straight wire), scaled by alpha
\draw[flow] (sp) -- (add);
\node[font=\sffamily\scriptsize, anchor=south] at (2.65, \yd+0.13)
  {{\color{dlMut}skip}\; $\times\,\alpha$};
% branch through inner Norm and f_j into the junction
\draw[flow, rounded corners=3pt] (sp) |- (inrm.west);
\draw[flow] (inrm.east) -- (fj.west);
\draw[flow, rounded corners=3pt] (fj.east) -| (add.south);
% then Norm, then output
\draw[flow] (add) -- (nrm);
\draw[flow] (nrm) -- (ro);

% beta annotation on f_j
\draw[draw=dlLine, line width=0.45pt, dash pattern=on 1.6pt off 1.4pt]
  (fj.south) -- ++(0, -0.22);
\node[font=\sffamily\scriptsize, anchor=north] at (3.10, 0.20)
  {{\color{dlMut}init gain}\, $\beta$};

% =========================================================================
%  (d) SCALING RULE
% =========================================================================
\node[head] at (10.63, 3.10) {{\bfseries (d)\; Scaling rule}};
\node[rectangle, rounded corners=2pt, draw=dlInk!75, fill=black!3,
      line width=0.8pt, minimum width=39mm, minimum height=27mm,
      text width=32mm, anchor=north east, align=left, inner sep=8pt,
      font=\sffamily\footnotesize]
  (rule) at (14.53, 2.70) {%
    {\bfseries DeepLoop} \; ($p = \tfrac{1}{2}$)\\[5pt]
    $\alpha = (2N)^{1/2}$\\[2pt]
    $\beta  = (8N)^{-1/2}$\\[5pt]
    % {\scriptsize\color{dlMut}accounts for weight tying:\\
     % each $\phi_j$ is visited $R$ times
     };

\end{tikzpicture}%
\endgroup
.
%  Requires: tikz (+ arrows.meta, calc libraries), amsmath.
%  Helvetica (phv, scaled 0.92) is scoped to this group only.
% ==========================================================================
\begingroup
\makeatletter
\def\Hv@scale{0.92}%
\makeatother
\renewcommand{\sfdefault}{phv}%
% ---- palette: two colorblind-safe accents (block identity) + neutrals ----
\definecolor{dlBlue}  {HTML}{4C72B0} % Block 1 / phi_1
\definecolor{dlOrange}{HTML}{DD8452} % Block 2 / phi_2
\definecolor{dlInk}   {HTML}{2B2F36} % main text
\definecolor{dlMut}   {HTML}{6A7280} % muted annotations
\definecolor{dlLine}  {HTML}{9AA0A8} % neutral borders
\definecolor{dlArrow} {HTML}{565D66} % arrows
\begin{tikzpicture}[
  font=\sffamily,
  >={Stealth[length=1.7mm,width=1.4mm]},
  line cap=round, line join=round,
  every node/.style={text=dlInk, outer sep=0pt},
  % small sublayer box (unrolled chain)
  subS/.style 2 args={rectangle, rounded corners=2pt, draw=#1,
      line width=0.6pt, fill=#2, minimum width=7.6mm, minimum height=6.2mm,
      inner sep=1pt, font=\sffamily\scriptsize},
  % large sublayer box (physical panel)
  subL/.style 2 args={rectangle, rounded corners=2pt, draw=#1,
      line width=0.6pt, fill=#2, minimum width=11.5mm, minimum height=6.8mm,
      inner sep=1pt, font=\sffamily\footnotesize},
  % state chips x_0, x_N, x_i, ...
  chip/.style={rectangle, rounded corners=2pt, draw=dlLine, fill=black!5,
      line width=0.6pt, minimum height=6.2mm, minimum width=8.5mm,
      inner sep=2.5pt, font=\footnotesize},
  % neutral gray box (Norm, f_j)
  gray box/.style={rectangle, rounded corners=2pt, draw=dlLine, fill=black!4,
      line width=0.6pt, font=\sffamily\footnotesize},
  flow/.style={->, draw=dlArrow, line width=0.6pt,
      shorten >=1.2pt, shorten <=1.2pt},
  head/.style={font=\sffamily\small, anchor=base west, inner sep=0pt},
  mut/.style={font=\sffamily\scriptsize, text=dlMut},
]

% =========================================================================
%  (a) PHYSICAL BLOCKS
% =========================================================================
\node[head] at (-1.53, 10.15)
  {{\bfseries (a)\; Physical blocks} \textcolor{dlMut}{---\;
   $K{=}2$ shared blocks, stored once}};

% Block 1
\node[rectangle, rounded corners=2pt, draw=dlBlue, fill=dlBlue!6,
      line width=0.8pt, minimum width=38mm, minimum height=16mm]
  (B1) at (0.77, 8.90) {};
\node[font=\sffamily\footnotesize\bfseries, text=dlBlue]
  at (0.77, 9.37) {Block 1\; ($\phi_1$)};
\node[subL={dlBlue}{dlBlue!13}, minimum width=13mm]
  (b1a) at (-0.08, 8.60) {attn};
\node[subL={dlBlue}{dlBlue!30}, minimum width=13mm]
  (b1f) at ( 1.62, 8.60) {ffn};
\draw[flow] (b1a) -- (b1f);

% Block 2
\node[rectangle, rounded corners=2pt, draw=dlOrange, fill=dlOrange!6,
      line width=0.8pt, minimum width=38mm, minimum height=16mm]
  (B2) at (5.37, 8.90) {};
\node[font=\sffamily\footnotesize\bfseries, text=dlOrange]
  at (5.37, 9.37) {Block 2\; ($\phi_2$)};
\node[subL={dlOrange}{dlOrange!13}, minimum width=13mm]
  (b2a) at (4.52, 8.60) {attn};
\node[subL={dlOrange}{dlOrange!30}, minimum width=13mm]
  (b2f) at (6.22, 8.60) {ffn};
\draw[flow] (b2a) -- (b2f);

\draw[flow] (B1) -- (B2);

% loop-back arc under the two blocks, label set on the wire
\draw[flow, rounded corners=3pt]
  (B2.east) -- ++(0.40,0) |- (3.07, 7.55) -| ($(B1.west)+(-0.40,0)$)
  -- (B1.west);
\node[font=\sffamily\footnotesize, fill=white, inner sep=2.5pt]
  at (3.07, 7.55) {loop $\times\,R$ rounds};

% =========================================================================
%  (b) UNROLLED EXECUTION
% =========================================================================
\node[head] at (-1.53, 6.45)
  {{\bfseries (b)\; Unrolled execution} \textcolor{dlMut}{---\;
   the same two blocks are revisited in every round}};

\def\yc{5.45}
% input / output chips
\node[chip] (xin)  at (-1.105, \yc) {$\mathbf{x}_0$};
\node[chip] (xout) at (14.085, \yc) {$\mathbf{x}_N$};

% six unrolled 2-sublayer groups
\foreach \gx/\hue/\pj [count=\i] in
  {0/dlBlue/1, 2.30/dlOrange/2, 4.82/dlBlue/1,
   7.12/dlOrange/2, 9.64/dlBlue/1, 11.94/dlOrange/2}{
  \node[subS={\hue}{\hue!13}] (a\i) at (\gx, \yc) {attn};
  \node[subS={\hue}{\hue!30}] (f\i) at (\gx+1.04, \yc) {ffn};
  \draw[flow] (a\i) -- (f\i);
  \draw[draw=\hue!45, line width=0.5pt, rounded corners=2pt]
    (\gx-0.51, \yc-0.45) rectangle (\gx+1.55, \yc+0.45);
  \node[font=\sffamily\scriptsize, text=\hue] at (\gx+0.52, \yc-0.72)
    {$\phi_{\pj}$};
}

% chain arrows
\draw[flow] (xin) -- (a1);
\foreach \i [evaluate=\i as \j using int(\i+1)] in {1,...,5}{
  \draw[flow] (f\i) -- (a\j);}
\draw[flow] (f6) -- (xout);

% round dividers + labels
\draw[dlLine!70, line width=0.45pt, dash pattern=on 2.6pt off 2pt]
  (4.08, \yc+0.5) -- (4.08, \yc-1.25);
\draw[dlLine!70, line width=0.45pt, dash pattern=on 2.6pt off 2pt]
  (8.90, \yc+0.5) -- (8.90, \yc-1.25);
\node[mut] at ( 1.67, \yc-1.05) {round $r{=}1$};
\node[mut] at ( 6.49, \yc-1.05) {round $r{=}2$};
\node[mut] at (11.31, \yc-1.05) {round $r{=}3$};

% summary line
\node[font=\sffamily\footnotesize, text=dlInk] at (6.50, 3.90)
  {$R{=}3$ rounds $\times$ $K{=}2$ blocks $\;\Rightarrow\;$
   unrolled depth $N{=}KR{=}6$,\quad $M{=}2N{=}12$ sublayer visits};

% =========================================================================
%  (c) RESIDUAL SUBLAYER
% =========================================================================
\node[head] at (-1.53, 3.10) {{\bfseries (c)\; Residual sublayer:}\;\;
  $\mathbf{x}_{i+1} = \mathrm{Norm}\bigl(\alpha\,\mathbf{x}_i
    + f_j(\mathrm{Norm}(\mathbf{x}_i);\phi_j)\bigr)$};

% dataflow: skip path is the straight wire, branch dips through f_j
\def\yd{1.95}
\node[chip] (ri) at (-1.105, \yd) {$\mathbf{x}_i$};
\node[circle, draw=dlArrow, fill=white, line width=0.6pt,
      minimum size=4.2mm, inner sep=0pt, font=\footnotesize]
  (add) at (5.05, \yd) {$+$};
\node[gray box, minimum width=13.5mm, minimum height=6.5mm]
  (nrm) at (6.90, \yd) {Norm};
\node[chip] (ro) at (8.80, \yd) {$\mathbf{x}_{i+1}$};
\node[gray box, minimum width=10mm, minimum height=6.5mm]
  (inrm) at (1.15, 0.85) {Norm};
\node[gray box, minimum width=21mm, minimum height=8mm]
  (fj) at (3.10, 0.85) {$f_j(\,\cdot\,;\phi_j)$};

% split point
\coordinate (sp) at (0.25, \yd);
\fill[dlArrow] (sp) circle (1.1pt);
\draw[draw=dlArrow, line width=0.6pt] (ri) -- (sp);
% skip path (straight wire), scaled by alpha
\draw[flow] (sp) -- (add);
\node[font=\sffamily\scriptsize, anchor=south] at (2.65, \yd+0.13)
  {{\color{dlMut}skip}\; $\times\,\alpha$};
% branch through inner Norm and f_j into the junction
\draw[flow, rounded corners=3pt] (sp) |- (inrm.west);
\draw[flow] (inrm.east) -- (fj.west);
\draw[flow, rounded corners=3pt] (fj.east) -| (add.south);
% then Norm, then output
\draw[flow] (add) -- (nrm);
\draw[flow] (nrm) -- (ro);

% beta annotation on f_j
\draw[draw=dlLine, line width=0.45pt, dash pattern=on 1.6pt off 1.4pt]
  (fj.south) -- ++(0, -0.22);
\node[font=\sffamily\scriptsize, anchor=north] at (3.10, 0.20)
  {{\color{dlMut}init gain}\, $\beta$};

% =========================================================================
%  (d) SCALING RULE
% =========================================================================
\node[head] at (10.63, 3.10) {{\bfseries (d)\; Scaling rule}};
\node[rectangle, rounded corners=2pt, draw=dlInk!75, fill=black!3,
      line width=0.8pt, minimum width=39mm, minimum height=27mm,
      text width=32mm, anchor=north east, align=left, inner sep=8pt,
      font=\sffamily\footnotesize]
  (rule) at (14.53, 2.70) {%
    {\bfseries DeepLoop} \; ($p = \tfrac{1}{2}$)\\[5pt]
    $\alpha = (2N)^{1/2}$\\[2pt]
    $\beta  = (8N)^{-1/2}$\\[5pt]
    % {\scriptsize\color{dlMut}accounts for weight tying:\\
     % each $\phi_j$ is visited $R$ times
     };

\end{tikzpicture}%
\endgroup
.
%  Requires: tikz (+ arrows.meta, calc libraries), amsmath.
%  Helvetica (phv, scaled 0.92) is scoped to this group only.
% ==========================================================================
\begingroup
\makeatletter
\def\Hv@scale{0.92}%
\makeatother
\renewcommand{\sfdefault}{phv}%
% ---- palette: two colorblind-safe accents (block identity) + neutrals ----
\definecolor{dlBlue}  {HTML}{4C72B0} % Block 1 / phi_1
\definecolor{dlOrange}{HTML}{DD8452} % Block 2 / phi_2
\definecolor{dlInk}   {HTML}{2B2F36} % main text
\definecolor{dlMut}   {HTML}{6A7280} % muted annotations
\definecolor{dlLine}  {HTML}{9AA0A8} % neutral borders
\definecolor{dlArrow} {HTML}{565D66} % arrows
\begin{tikzpicture}[
  font=\sffamily,
  >={Stealth[length=1.7mm,width=1.4mm]},
  line cap=round, line join=round,
  every node/.style={text=dlInk, outer sep=0pt},
  % small sublayer box (unrolled chain)
  subS/.style 2 args={rectangle, rounded corners=2pt, draw=#1,
      line width=0.6pt, fill=#2, minimum width=7.6mm, minimum height=6.2mm,
      inner sep=1pt, font=\sffamily\scriptsize},
  % large sublayer box (physical panel)
  subL/.style 2 args={rectangle, rounded corners=2pt, draw=#1,
      line width=0.6pt, fill=#2, minimum width=11.5mm, minimum height=6.8mm,
      inner sep=1pt, font=\sffamily\footnotesize},
  % state chips x_0, x_N, x_i, ...
  chip/.style={rectangle, rounded corners=2pt, draw=dlLine, fill=black!5,
      line width=0.6pt, minimum height=6.2mm, minimum width=8.5mm,
      inner sep=2.5pt, font=\footnotesize},
  % neutral gray box (Norm, f_j)
  gray box/.style={rectangle, rounded corners=2pt, draw=dlLine, fill=black!4,
      line width=0.6pt, font=\sffamily\footnotesize},
  flow/.style={->, draw=dlArrow, line width=0.6pt,
      shorten >=1.2pt, shorten <=1.2pt},
  head/.style={font=\sffamily\small, anchor=base west, inner sep=0pt},
  mut/.style={font=\sffamily\scriptsize, text=dlMut},
]

% =========================================================================
%  (a) PHYSICAL BLOCKS
% =========================================================================
\node[head] at (-1.53, 10.15)
  {{\bfseries (a)\; Physical blocks} \textcolor{dlMut}{---\;
   $K{=}2$ shared blocks, stored once}};

% Block 1
\node[rectangle, rounded corners=2pt, draw=dlBlue, fill=dlBlue!6,
      line width=0.8pt, minimum width=38mm, minimum height=16mm]
  (B1) at (0.77, 8.90) {};
\node[font=\sffamily\footnotesize\bfseries, text=dlBlue]
  at (0.77, 9.37) {Block 1\; ($\phi_1$)};
\node[subL={dlBlue}{dlBlue!13}, minimum width=13mm]
  (b1a) at (-0.08, 8.60) {attn};
\node[subL={dlBlue}{dlBlue!30}, minimum width=13mm]
  (b1f) at ( 1.62, 8.60) {ffn};
\draw[flow] (b1a) -- (b1f);

% Block 2
\node[rectangle, rounded corners=2pt, draw=dlOrange, fill=dlOrange!6,
      line width=0.8pt, minimum width=38mm, minimum height=16mm]
  (B2) at (5.37, 8.90) {};
\node[font=\sffamily\footnotesize\bfseries, text=dlOrange]
  at (5.37, 9.37) {Block 2\; ($\phi_2$)};
\node[subL={dlOrange}{dlOrange!13}, minimum width=13mm]
  (b2a) at (4.52, 8.60) {attn};
\node[subL={dlOrange}{dlOrange!30}, minimum width=13mm]
  (b2f) at (6.22, 8.60) {ffn};
\draw[flow] (b2a) -- (b2f);

\draw[flow] (B1) -- (B2);

% loop-back arc under the two blocks, label set on the wire
\draw[flow, rounded corners=3pt]
  (B2.east) -- ++(0.40,0) |- (3.07, 7.55) -| ($(B1.west)+(-0.40,0)$)
  -- (B1.west);
\node[font=\sffamily\footnotesize, fill=white, inner sep=2.5pt]
  at (3.07, 7.55) {loop $\times\,R$ rounds};

% =========================================================================
%  (b) UNROLLED EXECUTION
% =========================================================================
\node[head] at (-1.53, 6.45)
  {{\bfseries (b)\; Unrolled execution} \textcolor{dlMut}{---\;
   the same two blocks are revisited in every round}};

\def\yc{5.45}
% input / output chips
\node[chip] (xin)  at (-1.105, \yc) {$\mathbf{x}_0$};
\node[chip] (xout) at (14.085, \yc) {$\mathbf{x}_N$};

% six unrolled 2-sublayer groups
\foreach \gx/\hue/\pj [count=\i] in
  {0/dlBlue/1, 2.30/dlOrange/2, 4.82/dlBlue/1,
   7.12/dlOrange/2, 9.64/dlBlue/1, 11.94/dlOrange/2}{
  \node[subS={\hue}{\hue!13}] (a\i) at (\gx, \yc) {attn};
  \node[subS={\hue}{\hue!30}] (f\i) at (\gx+1.04, \yc) {ffn};
  \draw[flow] (a\i) -- (f\i);
  \draw[draw=\hue!45, line width=0.5pt, rounded corners=2pt]
    (\gx-0.51, \yc-0.45) rectangle (\gx+1.55, \yc+0.45);
  \node[font=\sffamily\scriptsize, text=\hue] at (\gx+0.52, \yc-0.72)
    {$\phi_{\pj}$};
}

% chain arrows
\draw[flow] (xin) -- (a1);
\foreach \i [evaluate=\i as \j using int(\i+1)] in {1,...,5}{
  \draw[flow] (f\i) -- (a\j);}
\draw[flow] (f6) -- (xout);

% round dividers + labels
\draw[dlLine!70, line width=0.45pt, dash pattern=on 2.6pt off 2pt]
  (4.08, \yc+0.5) -- (4.08, \yc-1.25);
\draw[dlLine!70, line width=0.45pt, dash pattern=on 2.6pt off 2pt]
  (8.90, \yc+0.5) -- (8.90, \yc-1.25);
\node[mut] at ( 1.67, \yc-1.05) {round $r{=}1$};
\node[mut] at ( 6.49, \yc-1.05) {round $r{=}2$};
\node[mut] at (11.31, \yc-1.05) {round $r{=}3$};

% summary line
\node[font=\sffamily\footnotesize, text=dlInk] at (6.50, 3.90)
  {$R{=}3$ rounds $\times$ $K{=}2$ blocks $\;\Rightarrow\;$
   unrolled depth $N{=}KR{=}6$,\quad $M{=}2N{=}12$ sublayer visits};

% =========================================================================
%  (c) RESIDUAL SUBLAYER
% =========================================================================
\node[head] at (-1.53, 3.10) {{\bfseries (c)\; Residual sublayer:}\;\;
  $\mathbf{x}_{i+1} = \mathrm{Norm}\bigl(\alpha\,\mathbf{x}_i
    + f_j(\mathrm{Norm}(\mathbf{x}_i);\phi_j)\bigr)$};

% dataflow: skip path is the straight wire, branch dips through f_j
\def\yd{1.95}
\node[chip] (ri) at (-1.105, \yd) {$\mathbf{x}_i$};
\node[circle, draw=dlArrow, fill=white, line width=0.6pt,
      minimum size=4.2mm, inner sep=0pt, font=\footnotesize]
  (add) at (5.05, \yd) {$+$};
\node[gray box, minimum width=13.5mm, minimum height=6.5mm]
  (nrm) at (6.90, \yd) {Norm};
\node[chip] (ro) at (8.80, \yd) {$\mathbf{x}_{i+1}$};
\node[gray box, minimum width=10mm, minimum height=6.5mm]
  (inrm) at (1.15, 0.85) {Norm};
\node[gray box, minimum width=21mm, minimum height=8mm]
  (fj) at (3.10, 0.85) {$f_j(\,\cdot\,;\phi_j)$};

% split point
\coordinate (sp) at (0.25, \yd);
\fill[dlArrow] (sp) circle (1.1pt);
\draw[draw=dlArrow, line width=0.6pt] (ri) -- (sp);
% skip path (straight wire), scaled by alpha
\draw[flow] (sp) -- (add);
\node[font=\sffamily\scriptsize, anchor=south] at (2.65, \yd+0.13)
  {{\color{dlMut}skip}\; $\times\,\alpha$};
% branch through inner Norm and f_j into the junction
\draw[flow, rounded corners=3pt] (sp) |- (inrm.west);
\draw[flow] (inrm.east) -- (fj.west);
\draw[flow, rounded corners=3pt] (fj.east) -| (add.south);
% then Norm, then output
\draw[flow] (add) -- (nrm);
\draw[flow] (nrm) -- (ro);

% beta annotation on f_j
\draw[draw=dlLine, line width=0.45pt, dash pattern=on 1.6pt off 1.4pt]
  (fj.south) -- ++(0, -0.22);
\node[font=\sffamily\scriptsize, anchor=north] at (3.10, 0.20)
  {{\color{dlMut}init gain}\, $\beta$};

% =========================================================================
%  (d) SCALING RULE
% =========================================================================
\node[head] at (10.63, 3.10) {{\bfseries (d)\; Scaling rule}};
\node[rectangle, rounded corners=2pt, draw=dlInk!75, fill=black!3,
      line width=0.8pt, minimum width=39mm, minimum height=27mm,
      text width=32mm, anchor=north east, align=left, inner sep=8pt,
      font=\sffamily\footnotesize]
  (rule) at (14.53, 2.70) {%
    {\bfseries DeepLoop} \; ($p = \tfrac{1}{2}$)\\[5pt]
    $\alpha = (2N)^{1/2}$\\[2pt]
    $\beta  = (8N)^{-1/2}$\\[5pt]
    % {\scriptsize\color{dlMut}accounts for weight tying:\\
     % each $\phi_j$ is visited $R$ times
     };

\end{tikzpicture}%
\endgroup
.
%  Requires: tikz (+ arrows.meta, calc libraries), amsmath.
%  Helvetica (phv, scaled 0.92) is scoped to this group only.
% ==========================================================================
\begingroup
\makeatletter
\def\Hv@scale{0.92}%
\makeatother
\renewcommand{\sfdefault}{phv}%
% ---- palette: two colorblind-safe accents (block identity) + neutrals ----
\definecolor{dlBlue}  {HTML}{4C72B0} % Block 1 / phi_1
\definecolor{dlOrange}{HTML}{DD8452} % Block 2 / phi_2
\definecolor{dlInk}   {HTML}{2B2F36} % main text
\definecolor{dlMut}   {HTML}{6A7280} % muted annotations
\definecolor{dlLine}  {HTML}{9AA0A8} % neutral borders
\definecolor{dlArrow} {HTML}{565D66} % arrows
\begin{tikzpicture}[
  font=\sffamily,
  >={Stealth[length=1.7mm,width=1.4mm]},
  line cap=round, line join=round,
  every node/.style={text=dlInk, outer sep=0pt},
  % small sublayer box (unrolled chain)
  subS/.style 2 args={rectangle, rounded corners=2pt, draw=#1,
      line width=0.6pt, fill=#2, minimum width=7.6mm, minimum height=6.2mm,
      inner sep=1pt, font=\sffamily\scriptsize},
  % large sublayer box (physical panel)
  subL/.style 2 args={rectangle, rounded corners=2pt, draw=#1,
      line width=0.6pt, fill=#2, minimum width=11.5mm, minimum height=6.8mm,
      inner sep=1pt, font=\sffamily\footnotesize},
  % state chips x_0, x_N, x_i, ...
  chip/.style={rectangle, rounded corners=2pt, draw=dlLine, fill=black!5,
      line width=0.6pt, minimum height=6.2mm, minimum width=8.5mm,
      inner sep=2.5pt, font=\footnotesize},
  % neutral gray box (Norm, f_j)
  gray box/.style={rectangle, rounded corners=2pt, draw=dlLine, fill=black!4,
      line width=0.6pt, font=\sffamily\footnotesize},
  flow/.style={->, draw=dlArrow, line width=0.6pt,
      shorten >=1.2pt, shorten <=1.2pt},
  head/.style={font=\sffamily\small, anchor=base west, inner sep=0pt},
  mut/.style={font=\sffamily\scriptsize, text=dlMut},
]

% =========================================================================
%  (a) PHYSICAL BLOCKS
% =========================================================================
\node[head] at (-1.53, 10.15)
  {{\bfseries (a)\; Physical blocks} \textcolor{dlMut}{---\;
   $K{=}2$ shared blocks, stored once}};

% Block 1
\node[rectangle, rounded corners=2pt, draw=dlBlue, fill=dlBlue!6,
      line width=0.8pt, minimum width=38mm, minimum height=16mm]
  (B1) at (0.77, 8.90) {};
\node[font=\sffamily\footnotesize\bfseries, text=dlBlue]
  at (0.77, 9.37) {Block 1\; ($\phi_1$)};
\node[subL={dlBlue}{dlBlue!13}, minimum width=13mm]
  (b1a) at (-0.08, 8.60) {attn};
\node[subL={dlBlue}{dlBlue!30}, minimum width=13mm]
  (b1f) at ( 1.62, 8.60) {ffn};
\draw[flow] (b1a) -- (b1f);

% Block 2
\node[rectangle, rounded corners=2pt, draw=dlOrange, fill=dlOrange!6,
      line width=0.8pt, minimum width=38mm, minimum height=16mm]
  (B2) at (5.37, 8.90) {};
\node[font=\sffamily\footnotesize\bfseries, text=dlOrange]
  at (5.37, 9.37) {Block 2\; ($\phi_2$)};
\node[subL={dlOrange}{dlOrange!13}, minimum width=13mm]
  (b2a) at (4.52, 8.60) {attn};
\node[subL={dlOrange}{dlOrange!30}, minimum width=13mm]
  (b2f) at (6.22, 8.60) {ffn};
\draw[flow] (b2a) -- (b2f);

\draw[flow] (B1) -- (B2);

% loop-back arc under the two blocks, label set on the wire
\draw[flow, rounded corners=3pt]
  (B2.east) -- ++(0.40,0) |- (3.07, 7.55) -| ($(B1.west)+(-0.40,0)$)
  -- (B1.west);
\node[font=\sffamily\footnotesize, fill=white, inner sep=2.5pt]
  at (3.07, 7.55) {loop $\times\,R$ rounds};

% =========================================================================
%  (b) UNROLLED EXECUTION
% =========================================================================
\node[head] at (-1.53, 6.45)
  {{\bfseries (b)\; Unrolled execution} \textcolor{dlMut}{---\;
   the same two blocks are revisited in every round}};

\def\yc{5.45}
% input / output chips
\node[chip] (xin)  at (-1.105, \yc) {$\mathbf{x}_0$};
\node[chip] (xout) at (14.085, \yc) {$\mathbf{x}_N$};

% six unrolled 2-sublayer groups
\foreach \gx/\hue/\pj [count=\i] in
  {0/dlBlue/1, 2.30/dlOrange/2, 4.82/dlBlue/1,
   7.12/dlOrange/2, 9.64/dlBlue/1, 11.94/dlOrange/2}{
  \node[subS={\hue}{\hue!13}] (a\i) at (\gx, \yc) {attn};
  \node[subS={\hue}{\hue!30}] (f\i) at (\gx+1.04, \yc) {ffn};
  \draw[flow] (a\i) -- (f\i);
  \draw[draw=\hue!45, line width=0.5pt, rounded corners=2pt]
    (\gx-0.51, \yc-0.45) rectangle (\gx+1.55, \yc+0.45);
  \node[font=\sffamily\scriptsize, text=\hue] at (\gx+0.52, \yc-0.72)
    {$\phi_{\pj}$};
}

% chain arrows
\draw[flow] (xin) -- (a1);
\foreach \i [evaluate=\i as \j using int(\i+1)] in {1,...,5}{
  \draw[flow] (f\i) -- (a\j);}
\draw[flow] (f6) -- (xout);

% round dividers + labels
\draw[dlLine!70, line width=0.45pt, dash pattern=on 2.6pt off 2pt]
  (4.08, \yc+0.5) -- (4.08, \yc-1.25);
\draw[dlLine!70, line width=0.45pt, dash pattern=on 2.6pt off 2pt]
  (8.90, \yc+0.5) -- (8.90, \yc-1.25);
\node[mut] at ( 1.67, \yc-1.05) {round $r{=}1$};
\node[mut] at ( 6.49, \yc-1.05) {round $r{=}2$};
\node[mut] at (11.31, \yc-1.05) {round $r{=}3$};

% summary line
\node[font=\sffamily\footnotesize, text=dlInk] at (6.50, 3.90)
  {$R{=}3$ rounds $\times$ $K{=}2$ blocks $\;\Rightarrow\;$
   unrolled depth $N{=}KR{=}6$,\quad $M{=}2N{=}12$ sublayer visits};

% =========================================================================
%  (c) RESIDUAL SUBLAYER
% =========================================================================
\node[head] at (-1.53, 3.10) {{\bfseries (c)\; Residual sublayer:}\;\;
  $\mathbf{x}_{i+1} = \mathrm{Norm}\bigl(\alpha\,\mathbf{x}_i
    + f_j(\mathrm{Norm}(\mathbf{x}_i);\phi_j)\bigr)$};

% dataflow: skip path is the straight wire, branch dips through f_j
\def\yd{1.95}
\node[chip] (ri) at (-1.105, \yd) {$\mathbf{x}_i$};
\node[circle, draw=dlArrow, fill=white, line width=0.6pt,
      minimum size=4.2mm, inner sep=0pt, font=\footnotesize]
  (add) at (5.05, \yd) {$+$};
\node[gray box, minimum width=13.5mm, minimum height=6.5mm]
  (nrm) at (6.90, \yd) {Norm};
\node[chip] (ro) at (8.80, \yd) {$\mathbf{x}_{i+1}$};
\node[gray box, minimum width=10mm, minimum height=6.5mm]
  (inrm) at (1.15, 0.85) {Norm};
\node[gray box, minimum width=21mm, minimum height=8mm]
  (fj) at (3.10, 0.85) {$f_j(\,\cdot\,;\phi_j)$};

% split point
\coordinate (sp) at (0.25, \yd);
\fill[dlArrow] (sp) circle (1.1pt);
\draw[draw=dlArrow, line width=0.6pt] (ri) -- (sp);
% skip path (straight wire), scaled by alpha
\draw[flow] (sp) -- (add);
\node[font=\sffamily\scriptsize, anchor=south] at (2.65, \yd+0.13)
  {{\color{dlMut}skip}\; $\times\,\alpha$};
% branch through inner Norm and f_j into the junction
\draw[flow, rounded corners=3pt] (sp) |- (inrm.west);
\draw[flow] (inrm.east) -- (fj.west);
\draw[flow, rounded corners=3pt] (fj.east) -| (add.south);
% then Norm, then output
\draw[flow] (add) -- (nrm);
\draw[flow] (nrm) -- (ro);

% beta annotation on f_j
\draw[draw=dlLine, line width=0.45pt, dash pattern=on 1.6pt off 1.4pt]
  (fj.south) -- ++(0, -0.22);
\node[font=\sffamily\scriptsize, anchor=north] at (3.10, 0.20)
  {{\color{dlMut}init gain}\, $\beta$};

% =========================================================================
%  (d) SCALING RULE
% =========================================================================
\node[head] at (10.63, 3.10) {{\bfseries (d)\; Scaling rule}};
\node[rectangle, rounded corners=2pt, draw=dlInk!75, fill=black!3,
      line width=0.8pt, minimum width=39mm, minimum height=27mm,
      text width=32mm, anchor=north east, align=left, inner sep=8pt,
      font=\sffamily\footnotesize]
  (rule) at (14.53, 2.70) {%
    {\bfseries DeepLoop} \; ($p = \tfrac{1}{2}$)\\[5pt]
    $\alpha = (2N)^{1/2}$\\[2pt]
    $\beta  = (8N)^{-1/2}$\\[5pt]
    % {\scriptsize\color{dlMut}accounts for weight tying:\\
     % each $\phi_j$ is visited $R$ times
     };

\end{tikzpicture}%
\endgroup
